\begin{center}\textbf{Resumen}\end{center}
\begingroup
\fontsize{9.5}{10.7}\selectfont
\setlength{\parskip}{1.5pt plus .5pt minus .5pt}
\noindent
Se construye una prolongación de la estructura discreta del continuo que
relaciona sus coordenadas numéricas \(\pi,\varphi,e,\alpha\) con la incidencia excepcional, las
álgebras de operadores de vértice y la dualidad de un sector compacto.
La Holografía Modular Triádica (HMT) parte de \emph{All Possible Paths}
(APP): el soporte \((\mathbb Z/9\mathbb Z)^2\), su grafo de Cayley aditivo
con generadores cardinales \(\pm(1,0),\pm(0,1)\) y las evaluaciones suma
y producto, con residuos y cocientes conservados. La firma
\(\mathrm{TRIT}=\{-1,0,+1\}\) distingue orientación y régimen; el
\emph{Topological Phase Kernel} (TPK) compone selección de evaluación,
transporte y actualización de memoria para prolongar estados compatibles.
El mismo árbol de historias sostiene la medida cilíndrica, el transporte
de hojas, el balance orientado, la incidencia y los complejos de memoria
con sus mapas de refinamiento.

El registro dodecafásico generado \(K\) enlaza dos realizaciones.
Su incidencia determina una bandera de Witt que, con el pegado, la métrica
y la cámara declarados, conduce a un vecino de Leech. Los teoremas de
construcción de álgebras de vértice aplicables a ese retículo identifican
las extensiones de órdenes dos y tres con \(V^\natural\), cuyo grupo de
automorfismos es el Monstruo; se transportan productos, graduación y
trazas mediante isomorfismos elegidos. En la representación de doce
posiciones de \(A_5\), la proyección de \(K\) sobre el sector
tridimensional es no nula y determina la reducción ortogonal
\(12\to11\to10\). La órbita cíclica del registro permite identificar
operadores con cotas de estabilidad; el balance bilateral proporciona
recuperación de estado y memoria bajo sus hipótesis isométricas.

% BEGIN S27_FRAMING
La comparación de las estrellas de \(B_K\), \(B_\varphi\) y \(B_{\rm APP}\) distingue
las firmas polarizadas \((3,1,0)\), \((1,0,3)\) y \((1,1,2)\).
El marco ordenado de los cuatro soportes
\(H^-_\alpha,H_e,H_\varphi,H_{\rm APP}\) tiene estabilizador trivial
en la acción de Witt sobre las doce posiciones. La lectura numérica
conserva, además de los residuos regionales, los cocientes enteros de
suma y diferencia: su composición satisface una identidad de cociclo,
y sus reducciones aportan componentes algebraicas adicionales a la
evaluación afín de las entradas. Estos resultados precisan las
funciones complementarias de incidencia, orientación y memoria de
acarreo dentro de la prolongación común.

% END S27_FRAMING
La realización de Weyl de la fase genera las álgebras matriciales de cada
profundidad; las inclusiones de fibra completa producen un álgebra
uniformemente hiperfinita y un factor tracial de tipo \(II_1\).
La sección de acción y el par angular fijan una elipse orientada y un
radio normalizado único. El intercambio de coordenadas enteras y la
inversión del radio conservan su forma cuadrática y entrelazan
unitariamente los Hamiltonianos autoadjuntos, incluidos sus dominios.
El transporte de la relación residual garantiza la covariancia del cambio
de sección. Un morfismo de pantallas coordina esa dualidad con la reducción
dirigida por \(K\) y el cociente isotrópico reticular \(26\to24\).

La composición de trayectorias con frontera conservada define un cociclo
aditivo de acción. Su prolongación hacia la teoría M especifica el
codominio de campos undecadimensionales y los mapas de realización que
deben conservar acción, graduación y dualidad. Bajo los entrelazamientos
y dominios declarados, la identidad de supercarga se transporta a la imagen
isométrica. Las condiciones de identificación física completa se
distinguen de estos resultados algebraicos y operatorios.
\par\medskip

% BEGIN R27_FRAMING_SUMMARY_ES
La aritmética que antecede a esas realizaciones incluye las leyes del
defecto APP bajo cambio de módulo y dimensión, la subretícula invariante
de índice dos y la codificación nonádica con acarreo. El emisor determina
un alfabeto de \(42\) bloques y empalmes regionales explícitos; el registro
ordenado conserva correlaciones que distingue de su histograma. La
correlación cuadrática demuestra la identidad de la coordenatización de frontera,
y la selección separa saturación, neutralidad dual y orientación. La
identidad racional de nivel tres establece además una divisibilidad
uniforme óptima por \(3^9\) en el anillo de series formales.
% END R27_FRAMING_SUMMARY_ES

\noindent\textbf{Palabras clave:} Holografía Modular Triádica; estructura
discreta del continuo; registro dodecafásico; álgebras de operadores de
vértice; Moonshine; Monstruo; dualidad T; representación de Weyl;
pantallas dimensionales; teoría M.
\par\endgroup
